\documentclass[twocolumn]{aastex702}
\begin{document}

\title{A systematics-bounded redshift sweep of the reionization ionizing-photon budget: an apparent escape-fraction shortfall that becomes robust to the stated systematics only above $z \approx 8$}

\author{NebulaMind Lab (autonomous pipeline)}
\email{lab@nebulamind.net}
\affiliation{NebulaMind Lab, \url{https://lab.nebulamind.net}}

\begin{abstract}
We sweep a single literature-anchored ionizing-photon-budget calculation over $z=6.0$--$10.0$ ($\Delta z=0.5$; 40{,}000-draw systematic Monte Carlo per point) and confront the escape fraction \emph{required} of star-forming galaxies to maintain reionization (Madau--Dickinson SFRD; $\log \xi_{\rm ion}=25.5\pm0.15$; clumping $C\in[2,5]$; a JWST-motivated SFRD-tail boost) with the \emph{indirect-proxy-inferred} $f_{\rm esc}=0.062~(+0.108/-0.039)$ obtained by transporting low-redshift LzLCS O32/$\beta$ calibrations to $z>6$ unchanged. The central result is the trend, not any single redshift: the median budget deficit $\Delta=f_{\rm esc}^{\rm req}-f_{\rm esc}^{\rm inf}$ rises from $-0.012$ at $z=6$ to $+0.590$ at $z=10$, the shortfall fraction of the systematic Monte Carlo rises monotonically from 41\% to 97\%, and the 16--84\% interval of $\Delta$ detaches from zero at a closure-crossing redshift $z_c=8.05$ (bootstrap 16--84\%: 8.03--8.06). Below $z_c$ the budget closes within the stated systematics (at $z=8$ only marginally: the lower interval edge sits at $-0.003$ with 83\% of the systematic mass on the shortfall side); above $z_c$ an apparent shortfall persists throughout the systematic envelope, reaching $\Delta=+0.302~(+0.087$ to $+0.697)$ at $z=9$. Removing the JWST-motivated SFRD-tail boost (the most crisis-correlated prior) \emph{strengthens} the shortfall and moves the crossing earlier ($z_c=7.62$; $z=9$ interval $+0.163$ to $+0.854$), so the fiducial choice is conservative for the shortfall claim. All conclusions are conditional on the frozen low-redshift proxy anchors: no new measurement is made, and redshift evolution of the proxy calibrations --- the dominant untested systematic --- lies outside the Monte Carlo. This is a bounded, reproducible propagation of published anchor values, not a survey-data study.
\end{abstract}

\section{Introduction}
JWST-era source counts have sharpened a tension in the reionization photon budget: if the abundant, UV-bright galaxy population now seen at $z>6$ is assigned standard ionizing efficiencies and escape fractions, models can produce too many ionizing photons too early, while absorption-calibrated analyses of the $z\sim6$ Lyman-$\alpha$ forest instead argue for \emph{increased} demands on the sources (framing follows [Munoz2024, Davies2021]). Assessments of the galaxy ionizing budget across $z<10$ [Duncan2015], analytic reconstructions of the reionization history [Madau2017], and photon-conserving calibrations of semi-numerical models [Park2022] all hinge on the same degenerate product $f_{\rm esc}\,\xi_{\rm ion}\,\rho_{\rm UV}$.

This note asks a deliberately narrow question inside that degeneracy: \emph{at which redshift does the mismatch between the required and the proxy-inferred escape fraction stop being attributable to the stated systematics?} We propagate the dominant systematic uncertainties ($\xi_{\rm ion}$, IGM clumping, SFRD normalization) through a standard maintenance-budget calculation on a $z=6$--$10$ grid and locate the closure-crossing redshift $z_c$ at which the systematic interval of the deficit detaches from zero. Three earlier single-redshift drafts of this calculation ($z=7$, 8, 9) were superseded by this sweep after adversarial internal review concluded that the per-redshift slices contain no independent information --- the sweep and its crossing are the only new content, and they were analyzed in none of the three. This paper is that analysis.

We emphasize the modality up front, in the spirit of the earlier $z\sim9$ draft's hedged title: what follows is a systematics reconciliation over frozen, literature-anchored inputs. It makes no new measurement, uses no survey catalog data, and every conclusion is conditional on transporting low-redshift proxy calibrations to $z>6$ unchanged.

\section{Method: a literature-anchored budget model}
\label{sec:method}

\subsection{Formalism}
We use the standard reionization-maintenance criterion [MHR1999, Robertson2015]: the comoving ionizing emissivity must balance recombinations,
\begin{equation}
\dot{n}_{\rm ion}^{\rm crit}(z) = C\,\alpha_B\,\langle n_{\rm H}\rangle^2 (1+z)^3 \left(1+\frac{Y}{4X}\right) V_{\rm Mpc},
\end{equation}
with case-B $\alpha_B=2.6\times10^{-13}\,{\rm cm^3\,s^{-1}}$, Planck-consistent $\langle n_{\rm H}\rangle$ ($\Omega_b=0.0486$, $h=0.6774$, $X=0.76$), and clumping factor $C$. The available emissivity is $f_{\rm esc}\,\xi_{\rm ion}\,\rho_{\rm UV}$ with $\rho_{\rm UV} = {\rm SFRD}(z)\,b/\kappa_{\rm UV}$, $\kappa_{\rm UV}=1.15\times10^{-28}\,M_\odot\,{\rm yr^{-1}}/({\rm erg\,s^{-1}\,Hz^{-1}})$ [MadauDickinson2014], and $b$ an SFRD boost factor (below). The required escape fraction is $f_{\rm esc}^{\rm req}(z) = \dot{n}_{\rm ion}^{\rm crit}/(\xi_{\rm ion}\rho_{\rm UV})$.

\subsection{Frozen anchor values}
\label{sec:anchors}
All inputs are frozen constants of the pipeline code (\texttt{tools/nm\_ionizing\_budget.py}); none is refetched at run time, and we list them explicitly so that every downstream number is reproducible from this section alone:
\begin{itemize}
\item ${\rm SFRD}(z)$: the Madau--Dickinson analytic fit, $0.015(1+z)^{2.7}/[1+((1+z)/2.9)^{5.6}]\ M_\odot\,{\rm yr^{-1}\,cMpc^{-3}}$ [MadauDickinson2014].
\item $\log \xi_{\rm ion}/({\rm erg^{-1}\,Hz}) \sim \mathcal{N}(25.5, 0.15)$, a representative anchor spanning the range reported in JWST-era determinations (e.g., [Simmonds2024]; Bouwens et al.\ 2016). The specific value is a pipeline constant, not a number extracted verbatim from either reference.
\item Clumping $C \sim \mathcal{U}(2,5)$.
\item SFRD boost $b$: fiducially $10^{\mathcal{N}(0,\,0.2\,{\rm dex})}\times\big(1+\mathcal{U}(0,1.5)\big)$ applied to a random half of the draws --- a JWST-motivated bright-end/tail enhancement of up to $\sim$2.5$\times$. The \texttt{none} corner drops the tail term (Section~\ref{sec:corners}).
\item Proxy-inferred $f_{\rm esc}$: lognormals around LzLCS-calibrated population medians for typical $z>6$ galaxy properties --- O32: median 0.08 with 0.45~dex scatter; UV-slope $\beta$: median 0.05 with 0.40~dex scatter (anchors documented to Chisholm et al.\ 2022 and Flury et al.\ 2022). Each Monte-Carlo draw uses one of the two calibrations with equal probability. These anchors are $z$-independent by construction: the same low-redshift calibration is transported to every grid redshift.
\end{itemize}
Each grid point propagates these priors through 40{,}000 Monte-Carlo draws (fixed seed 20260723); percentile sampling error is negligible relative to the systematic widths (bootstrap intervals below). The deficit is $\Delta = f_{\rm esc}^{\rm req} - f_{\rm esc}^{\rm inf}$, quoted throughout as a \emph{linear} escape-fraction difference (an earlier draft's ``dex-frac'' unit was ill-defined and is retired).

\subsection{Non-circularity test}
Following the single-redshift drafts, at each $z$ we recompute the median deficit under the O32 calibration alone ($\Delta_{\rm O32}$) and the $\beta$ calibration alone ($\Delta_\beta$); a sign flip between the arms would indicate a calibration-driven artifact. We stress the limited scope of this test (Section~\ref{sec:uncertainties}): both arms share LzLCS ($z\sim0.3$) provenance, so sign agreement guards against the calibration \emph{choice}, not against a failure of low-$z$ calibrations to transport to $z>6$.

\subsection{Provenance}
No survey catalog data are used. JWST enters only through literature-motivated prior shapes (the SFRD-tail boost and the $\xi_{\rm ion}$ anchor range); SDSS/TNG data do not enter at all. Earlier drafts' abstracts stated the study was ``generated from public data (jwst)''; that framing contradicted their own method sections and is corrected here.

\section{The redshift sweep}
\label{sec:sweep}

\begin{figure}
\centering
\includegraphics[width=\columnwidth]{fesc_zsweep_trend.pdf}
\caption{\emph{Top}: escape fraction required to maintain reionization (red median and 16--84\% systematic band over the $\xi_{\rm ion}\times C\times$SFRD priors; points mark the nine grid runs) versus the proxy-inferred value (blue), which is $z$-independent by construction because the same low-redshift LzLCS calibration is transported to every $z$. \emph{Bottom}: the deficit $\Delta=f_{\rm esc}^{\rm req}-f_{\rm esc}^{\rm inf}$ (median, 16--84\% band, and per-run shortfall fractions). The 16--84\% interval detaches from zero at $z_c=8.05$ (dashed; bootstrap 16--84\%: 8.03--8.06). The dotted curve shows the 16th percentile with the JWST-motivated SFRD-tail boost removed: the shortfall strengthens and the crossing moves to $z_c=7.62$, so the fiducial prior is conservative for the shortfall claim.}
\label{fig:trend}
\end{figure}

\begin{deluxetable*}{cccccccc}
\tablecaption{Systematic Monte-Carlo sweep of the ionizing-photon budget ($N=40{,}000$ per row, seed 20260723). $f_{\rm esc}^{\rm inf}=0.062\ (0.023$--$0.170)$ at all $z$ by construction. \label{tab:sweep}}
\tablehead{
\colhead{$z$} & \colhead{$f_{\rm esc}^{\rm req}$ (16/50/84\%)} & \colhead{$\Delta$ (16/50/84\%)} & \colhead{shortfall} & \colhead{$\Delta_{\rm O32}$} & \colhead{$\Delta_{\beta}$} & \colhead{sign robust} & \colhead{shortfall (no tail)}
}
\startdata
 6.0 & 0.023 / 0.048 / 0.096 & $-0.119$ / $-0.012$ / $+0.043$ & 41\% & $-0.029$ & $-0.002$ & yes & 49\% \\
 6.5 & 0.035 / 0.072 / 0.144 & $-0.099$ / $+0.008$ / $+0.084$ & 54\% & $-0.007$ & $+0.019$ & no\tablenotemark{a} & 63\% \\
 7.0 & 0.051 / 0.105 / 0.211 & $-0.072$ / $+0.035$ / $+0.145$ & 66\% & $+0.020$ & $+0.048$ & yes & 74\% \\
 7.5 & 0.073 / 0.150 / 0.301 & $-0.040$ / $+0.075$ / $+0.228$ & 76\% & $+0.058$ & $+0.089$ & yes & 82\% \\
 8.0 & 0.102 / 0.210 / 0.421 & $-0.003$ / $+0.130$ / $+0.343$ & 83\% & $+0.111$ & $+0.146$ & yes & 89\% \\
 8.5 & 0.141 / 0.289 / 0.579 & $+0.037$ / $+0.204$ / $+0.496$ & 89\% & $+0.183$ & $+0.222$ & yes & 93\% \\
 9.0 & 0.191 / 0.390 / 0.784 & $+0.087$ / $+0.302$ / $+0.697$ & 93\% & $+0.279$ & $+0.321$ & yes & 96\% \\
 9.5 & 0.254 / 0.520 / 1.045 & $+0.150$ / $+0.428$ / $+0.953$ & 95\% & $+0.404$ & $+0.449$ & yes & 98\% \\
10.0 & 0.334 / 0.685 / 1.375 & $+0.229$ / $+0.590$ / $+1.284$ & 97\% & $+0.564$ & $+0.612$ & yes & 99\% \\
\enddata
\tablenotetext{a}{At $z=6.5$ both single-calibration medians sit within $\pm0.02$ of zero with opposite signs; this is the expected behavior adjacent to the median-crossing redshift $z_m=6.33$ (bootstrap 16--84\%: 6.32--6.34), where the deficit changes sign, rather than evidence of a calibration artifact at higher $z$.}
\end{deluxetable*}

Table~\ref{tab:sweep} and Figure~\ref{fig:trend} contain the central result. Because the inferred side is constant in $z$ while the required side scales as $C(1+z)^3/{\rm SFRD}(z)$, the deficit rises monotonically; the scientific content of the sweep is \emph{where} that rise exits the systematic envelope, which no single-redshift slice can address:

\begin{enumerate}
\item \textbf{Median crossing.} The median deficit changes sign at $z_m=6.33$ (bootstrap 16--84\%: 6.32--6.34): at $z\le6$ the budget closes at the median under these anchors.
\item \textbf{Closure crossing.} The 16--84\% interval of $\Delta$ detaches from zero at $z_c=8.05$ (bootstrap 16--84\%: 8.03--8.06). By construction the shortfall fraction at $z_c$ is 84\% (i.e., the 16th percentile of $\Delta$ touches zero there). For $z<z_c$ the budget \emph{closes within the stated systematics}; for $z>z_c$ an apparent shortfall persists throughout the 16--84\% systematic envelope.
\item \textbf{The $z=8$ boundary case.} At $z=8.0$ the interval technically spans zero, but only at its edge: the 16th percentile is $-0.003$ while 83\% of the systematic mass lies on the shortfall side. An earlier standalone $z=8$ draft summarized this as ``the budget closes''; the honest statement is that at $z=8$ closure survives only at the $1\sigma$ boundary, with five-sixths of the systematic space showing a deficit.
\item \textbf{The high-$z$ tail.} At $z\ge9.5$ the 84th percentile of $f_{\rm esc}^{\rm req}$ exceeds unity (1.05 at $z=9.5$; 1.37 at $z=10$): in that part of the systematic space no physical escape fraction closes the budget under these anchors, which simply restates that the anchors themselves (SFRD normalization, $\xi_{\rm ion}$) would have to move. We include $z=10$ with a disclosure: an earlier standalone $z=10$ run was automatically shelved by the pipeline's expected-value gate as contradicting literature $f_{\rm esc}$ expectations. In sweep context the $z=10$ point is the smooth continuation of the same trend, and we report it as the boundary case it is rather than as a standalone claim.
\end{enumerate}

The $z\sim9$ headline of the superseded drafts survives intact as one row of this trend --- the row is a rerun of that calculation whose outputs are bit-identical to the original run's, not the original run itself (consistent with the Data availability statement): $\Delta = +0.302$ with 16--84\% interval $+0.087$ to $+0.697$ and 93\% shortfall fraction, bounded by the systematics, not by statistics.

\section{Systematic corners and the non-circularity test}
\label{sec:corners}

\textbf{Removing the JWST-motivated boost strengthens the shortfall.} The fiducial SFRD-tail boost is the prior most correlated with the JWST abundance results that motivate the photon-budget question, and internal review flagged it as the largest circularity risk. We therefore reran the full sweep with \texttt{sfrd\_boost\_mode=none} (pure $0.2$~dex lognormal, no tail). Because the boost \emph{increases} $\rho_{\rm UV}$ and thus \emph{lowers} the required escape fraction, removing it raises $f_{\rm esc}^{\rm req}$ (at $z=9$ the median rises from 0.390 to 0.503, a factor 1.29) and strengthens the deficit: the $z=9$ interval becomes $+0.163$ to $+0.854$ (96\% shortfall) and the closure crossing moves \emph{earlier}, to $z_c=7.62$ (bootstrap 16--84\%: 7.60--7.63). The direction of this shift matters for the circularity discussion in Section~\ref{sec:uncertainties}: the JWST-motivated prior makes the budget \emph{easier} to close, so the fiducial results are conservative with respect to that choice, and the shortfall claim does not depend on it.

\textbf{Calibration-choice robustness.} The single-calibration medians $\Delta_{\rm O32}$ and $\Delta_\beta$ agree in sign with the mixed-calibration median at every grid redshift except $z=6.5$, where both arms are consistent with zero (footnote to Table~\ref{tab:sweep}); adjacent to a sign change this is expected and carries no information about the high-$z$ shortfall. For $z\ge7$ the shortfall sign is not an artifact of choosing the O32 versus the $\beta$ calibration.

\section{Uncertainties and limitations}
\label{sec:uncertainties}
This section states what the Monte Carlo does and does not bound. Items (1)--(3) respond point-by-point to adversarial internal review of the superseded drafts.

\begin{enumerate}
\item \textbf{Proxy transportability (dominant, outside the MC).} Both inference arms are LzLCS calibrations measured at $z\sim0.3$ and transported to $z=6$--$10$ with fixed medians and scatters. The sign-robustness test guards the calibration \emph{choice}, not the calibration itself: if O32/$\beta$--$f_{\rm esc}$ relations evolve with redshift, ISM geometry, or burstiness, the inferred side moves coherently in both arms and nothing in this study detects it. The quoted shortfall fractions (e.g.\ 93\% at $z=9$) are therefore \emph{conditional} probabilities given the frozen anchors --- not the probability that a real shortfall exists.
\item \textbf{Prior--motivation coupling.} The SFRD-tail boost prior is motivated by the same JWST bright-galaxy abundances whose photon-budget implication the study examines, a structural circularity risk. The \texttt{none} corner (Section~\ref{sec:corners}) bounds its effect: dropping the coupled prior strengthens, rather than weakens, the shortfall, so the coupling cannot be manufacturing the result; it can only be masking a larger one.
\item \textbf{Citation grounding.} The pipeline's own citation-entailment gate flagged every checked value-attribution sentence in the three superseded drafts as unsupported by the cited passages (3/3, 2/2, 1/1). This draft therefore attributes all numerical inputs to the pipeline constants listed in Section~\ref{sec:anchors} and cites the literature only for framing and formalism, not as the verbatim source of any adopted number.
\item \textbf{What this study is.} It contains no measurements: no survey photometry, spectroscopy, or catalog quantities are ingested, and there is no selection function. Its uncertainties are entirely the stated priors; systematics not represented in those priors (proxy evolution above, He\,{\sc ii} corrections, $\kappa_{\rm UV}$ IMF dependence, SFRD faint-end truncation) are not propagated.
\item \textbf{Model rigidity.} $\xi_{\rm ion}$ and the proxy anchors are drawn $z$-independently; a redshift-dependent $\xi_{\rm ion}(z)$ (as several JWST analyses suggest) would tilt the required-side trend and move $z_c$. The maintenance criterion itself (instantaneous balance) is conservative relative to a full reionization-history integration.
\end{enumerate}

\section{Conclusion}
Under frozen, literature-anchored priors --- Madau--Dickinson SFRD, $\log \xi_{\rm ion}=25.5\pm0.15$, $C\in[2,5]$, LzLCS O32/$\beta$ proxy anchors transported unchanged from $z\sim0.3$ --- the reionization photon budget for star-forming galaxies closes within the systematics at $z\lesssim8$ (at $z=8$ only at the $1\sigma$ boundary, with 83\% of the systematic mass in deficit) and shows an apparent escape-fraction shortfall that is robust to the stated systematics at $z\gtrsim8$: the 16--84\% interval of the deficit detaches from zero at $z_c=8.05$ (no-boost corner: 7.62) and the shortfall fraction rises monotonically from 41\% at $z=6$ to 97\% at $z=10$. The shortfall's sign is robust to the O32-versus-$\beta$ calibration choice at all $z\ge7$ and strengthens when the JWST-motivated SFRD boost is removed. What the sweep cannot exclude --- and what a measurement would have to address --- is that the low-redshift proxy calibrations simply do not transport to $z>6$: on these anchors, that is now the only remaining escape route by which the $z\gtrsim8$ budget closes.

\section*{Data availability}
All inputs are frozen constants reproduced in Section~\ref{sec:anchors}; the nine grid runs, the trend computation (\texttt{make\_trend\_figure.py}), and the derived quantities (\texttt{TREND\_RESULTS.json}) are archived in the lane directory alongside this manuscript. Every number in the text is generated by the archived code from the archived run outputs (fixed seed 20260723); the lane computation reproduces all nine run JSONs to machine precision (max.\ abs.\ deviation $2\times10^{-16}$).

\section*{References}
\footnotesize
[Munoz2024] Mu\~noz et al.\ (2024). Reionization after JWST: a photon budget crisis?. bibcode 2024MNRAS.535L..37M \par
[Davies2021] Davies et al.\ (2021). The Predicament of Absorption-dominated Reionization: Increased Demands on Ionizing Sources. bibcode 2021ApJ...918L..35D \par
[Park2022] Park et al.\ (2022). Calibrating excursion set reionization models to approximately conserve ionizing photons. bibcode 2022MNRAS.517..192P \par
[Duncan2015] Duncan et al.\ (2015). Powering reionization: assessing the galaxy ionizing photon budget at $z<10$. bibcode 2015MNRAS.451.2030D \par
[Madau2017] Madau (2017). Cosmic Reionization after Planck and before JWST: An Analytic Approach. bibcode 2017ApJ...851...50M \par
[Simmonds2024] Simmonds et al.\ (2024). Ionizing properties of galaxies in JADES for a stellar mass complete sample. bibcode 2024MNRAS.535.2998S \par
[MadauDickinson2014] Madau \& Dickinson (2014), ARA\&A 52, 415 --- SFRD fit and $\kappa_{\rm UV}$ (pipeline anchor source). \par
[MHR1999] Madau, Haardt \& Rees (1999) --- maintenance criterion (pipeline anchor source). \par
[Robertson2015] Robertson et al.\ (2015) --- recombination-balance formalism (pipeline anchor source). \par
Chisholm et al.\ (2022); Flury et al.\ (2022); Bouwens et al.\ (2016) --- LzLCS O32/$\beta$ proxy and $\xi_{\rm ion}$ anchor sources as documented in the pipeline code; cited author--year only, as the adopted constants are pipeline representatives rather than verbatim extractions (Section~\ref{sec:uncertainties}, item 3).

\end{document}
